\section{Analytic classes and changes of physical coordinates}\label{sec:model}
The geometric constructions prescribe a smooth velocity first. This section
records exactly which other solutions are excluded by the energy argument.
Pressure assumptions are kept separate: continuity in a pressure norm, a
norm at each individual time, and an integrable pressure gradient are
different conditions.

Write $\mathcal F_\nu[U]=U_t+(U\cdot\nabla)U-\nu\Delta U$.
All noncompact results below use $\mathbb R^3$; their velocities and forces
have one fixed compact spatial support unless the individual statement
says otherwise. This property concerns the constructed solution, not a
competing solution. On $[0,T]$, write $CH^k=C([0,T];H^k(\mathbb R^3))$
and $C^1L^2=C^1([0,T];L^2(\mathbb R^3))$. Spatial constants in pressure
may depend on time.

\begin{lemma}[The pressure term without a pressure norm]\label{lem:b-gradient}
If $w,g\in L^2(\mathbb R^3;\mathbb R^3)$ satisfy
$\operatorname{div}w=0$ and $\operatorname{curl}g=0$ in distributions, then
$\langle w,g\rangle_{L^2}=0$. In particular, any distributional gradient
$g=\nabla p\in L^2$ has zero pairing with $w$, without an assumption that
$p\in L^2$.
\end{lemma}
\begin{proof}
An $L^2$ field is a tempered distribution. The differential identities,
initially tested on compactly supported smooth functions, extend to
Schwartz tests by multiplying by cutoffs tending to one and using $L^2$
convergence of the tests and their first derivatives. Fourier transformation
gives $\xi\cdot\widehat w=0$ and
$\xi_j\widehat g_k-\xi_k\widehat g_j=0$. These are identities of locally
integrable functions, so they hold almost everywhere. Away from $\xi=0$,
$\widehat g$ is parallel to $\xi$ and $\widehat w$ is perpendicular to it.
Their Hermitian product is zero. The exceptional point has measure zero;
Plancherel proves the assertion.
\end{proof}

\begin{lemma}[Residual realization and separate comparison classes]
\label{lem:b-comparison}\label{lem:p2-realization}
Let $\nu>0$ and let $U$ be a prescribed smooth divergence-free velocity
on $[0,\infty)\times\mathbb R^3$, with $U(0)=0$. On each finite interval,
assume $U\in CH^2\cap C^1L^2$ and $U,\nabla U$ are bounded. Put
$f=\mathcal F_\nu[U]$. Then $(U,0)$ is a solution. Its velocity is unique
in each of the following separately stated comparison classes, with the
same force and initial data:
\begin{enumerate}[label=(\roman*),itemsep=2pt]
\item Smooth classical $v\in CH^2\cap C^1L^2$, with bounded $v$, and a
pressure representative $p\in CH^1$.
\item Smooth classical $v\in CH^2\cap C^1L^2$, with bounded $v,\nabla v$,
and an $H^1$ pressure representative at each time, with no required time
continuity in that pressure norm.
\item Strong solutions with $v\in CH^4\cap C^1H^2$ and distributional
$\nabla p\in CL^2$.
\item Smooth classical $v\in CH^2\cap C^1L^2$, with bounded $v,\nabla v$,
and a pressure representative $p\in CL^2$.
\item Smooth classical $v\in CH^2\cap C^1L^2$, with bounded $v,\nabla v$,
without an additional integrability or normalization condition on pressure.
\end{enumerate}
The pressure gradient is unique. An $L^2$ pressure representative, when
specified, is zero; otherwise pressure is determined up to a function of
time. The reference material flow is a smooth diffeomorphism at every
finite time and exists for every finite forward time.

On a flat torus $\T_L^3$, the same conclusion holds in the classical class
where $v,v_t$, spatial derivatives through order two, and $p,\nabla p$ are
continuous on finite closed cylinders, with periodic mean-zero pressure.
No Sobolev conditions at infinity are involved in that assertion.
\end{lemma}
\begin{proof}
Substitution proves existence of the prescribed solution. Set $w=v-U$.
The difference equation is
\[
 w_t+(v\cdot\nabla)w+(w\cdot\nabla)U
    =\nu\Delta w-\nabla p,\qquad \operatorname{div}w=0.
\]
The stated velocity regularity implies $w\in C^1L^2\cap CH^2$, so
$\frac d{dt}\|w\|_2^2=2\langle w,w_t\rangle$ at each time. In case
(iii), the stronger assumptions imply these inclusions. They also imply
bounded velocity and spatial derivatives through order two: for $j\le2$,
Cauchy--Schwarz in Fourier space uses the finite integral
$\int_{\mathbb R^3}|\xi|^{2j}(1+|\xi|^2)^{-4}\,d\xi$.

For the transport term, insert a cutoff $\chi_R$ that equals one on the
ball of radius $R$, is supported in the ball of radius $2R$, and satisfies
$|\nabla\chi_R|\le C/R$. Its boundary error is at most
$CR^{-1}\|v\|_\infty\|w\|_2^2$, which tends to zero. Diffusion is
integrated by parts in $H^2$; alternatively its cutoff error is bounded by
$CR^{-1}\|w\|_2\|\nabla w\|_2$. The deformation term is bounded by
$\|\nabla U\|_{\infty,\mathrm{op}}\|w\|_2^2$.

We justify the pressure pairing according to the stated alternatives.
For (i) and (ii), at each fixed time approximate the chosen $H^1$ pressure
by compactly supported smooth functions; distributional incompressibility
and $w\in L^2$ give $\langle\nabla p,w\rangle=0$. Thus no continuity
of $p$ in time is used in (ii). For (iv), integration by parts on the
cutoff gives a pressure boundary term bounded by
$CR^{-1}\|p\|_2\|w\|_2$, also tending to zero. For (iii), apply
Lemma~\ref{lem:b-gradient} directly. Finally, in (v) the equation itself
gives at each time
\[
 \nabla p=f-v_t-(v\cdot\nabla)v+\nu\Delta v\in L^2,
 \qquad \|(v\cdot\nabla)v\|_2\le\|v\|_\infty\|\nabla v\|_2.
\]
Here $f\in L^2$ follows from the identical estimate for $U$ and its stated
time and space regularity. The smooth pressure gradient is curl free, so
Lemma~\ref{lem:b-gradient} again cancels the term. This does not choose an
$L^2$ representative of the pressure.

In every case the resulting pointwise energy inequality is
\[
 \frac12\frac d{dt}\|w\|_2^2+\nu\|\nabla w\|_2^2
 \le \|\nabla U\|_{\infty,\mathrm{op}}\|w\|_2^2.
\]
The coefficient is bounded on each finite interval and $w(0)=0$, so an
integrating factor gives $w=0$. The pressure has already disappeared before
time integration; the proof imposes no unlisted temporal pressure norm.
Substitution then gives $\nabla p=0$. A spatially constant $L^2$ function
on $\mathbb R^3$ is zero, which proves the normalization assertion.

On the torus every integration by parts is periodic and has no boundary
term. The specified classical regularity justifies norm differentiation,
diffusion, and pressure cancellation. The same energy inequality gives
uniqueness and the mean-zero normalization fixes pressure. In either domain,
the smooth reference velocity and its bounded spatial derivative give
unique material trajectories; bounded speed prevents finite-time escape.
Solving backwards on a finite interval gives the inverse flow, and smooth
ODE dependence gives its smoothness.
\end{proof}

Cases (i)--(v) identify the precise hypotheses used by the applications;
their presence in a common lemma does not replace a theorem's declared
comparison class by a different one. In particular, a pressure norm is not
silently imposed in a gradient-only application. The proof is the classical
energy comparison for prescribed smooth solutions \cite[Section~18]{Leray1934},
with its noncompact pressure pairings made explicit.

For the flow $\Phi_t$, differentiation of its trajectory equation gives
\[
 U(t)=(\partial_t\Phi_t)\circ\Phi_t^{-1},\qquad
 \partial_{tt}\Phi_t(a)=
 (f-\nabla p+\nu\Delta U)(t,\Phi_t(a)).
\]
These identities describe the material form of the same constructed solution;
they do not add an independent existence claim.

\begin{lemma}[Physical scaling]\label{lem:p2-chart}
For $a,b>0$, $x=x_0+ay$, and $s=bt$, let
$U(t,x)=abV(bt,(x-x_0)/a)$. Then trajectories are transformed by this change
of variables, incompressibility is preserved, and at the prescribed physical
viscosity $\nu$,
\[
 \mathcal F_\nu[U](t,x)=ab^2\left[
 V_s+(V\cdot\nabla_y)V-\frac{\nu}{a^2b}\Delta_yV
 \right](bt,(x-x_0)/a).
\]
\end{lemma}
\begin{proof}
The trajectory equation gives $\dot x=abV$; differentiating in space gives
$\operatorname{div}_xU=b\operatorname{div}_yV$. The three residual terms
scale by $ab^2,ab^2$, and $b/a$, respectively, proving the formula.
\end{proof}
Thus a side-$10$ torus is kept as a physical side-$10$ torus. If a chart
is rescaled, the residual is recomputed with the fixed $\nu$; the equation
is not identified with a differently normalized viscosity.

\begin{proposition}[Projection on a flat torus]\label{prop:projection-l2}
Let $g$ be a smooth force on $[0,\infty)\times\T_L^3$ with zero
spatial mean at each time. There is a unique smooth mean-zero potential
$\phi$ satisfying $\Delta\phi=\operatorname{div}g$. Put
$Qg=\nabla\phi$ and $Pg=g-Qg$. Then $Pg$ is solenoidal and mean zero,
and replacing $(g,p)$ by $(Pg,p-\phi)$ preserves the velocity and the
pressure normalization.

Use normalized spatial measure, frequencies $\xi_k=2\pi k/L$, and norms
\[
 \|G\|_{H^s}^2=\sum_{k\in\Z^3}(1+|\xi_k|^2)^s|\widehat G(k)|^2,
 \qquad
 \|G\|_{C^r}=\max_{|\beta|\le r}\|\partial_x^\beta G\|_\infty.
\]
For all integers $h,r\ge0$, at each time,
\begin{equation}\label{eq:projection-estimate}
 \begin{aligned}
 \max\{\|\partial_t^h Pg\|_{C^r},\|\partial_t^h Qg\|_{C^r}\}
     &\le S_L\|\partial_t^h g\|_{H^{r+2}},\\
 \|\partial_t^h\phi\|_{C^{r+1}}
     &\le S_Ld_L\|\partial_t^h g\|_{H^{r+2}},
 \end{aligned}
\end{equation}
where $S_L=(1+52(L/(2\pi))^4)^{1/2}$ and
$d_L=(1+(L/(2\pi))^2)^{1/2}$.

Uniform bounds, temporal $L^2$ bounds on spatial suprema, and qualitative
uniform decay of all mixed derivatives pass from $g$ to $Pg,Qg,\phi$.
A quantitative envelope for time order $h$ and output spatial order $r$
passes to $Pg,Qg$ through order $r$ and to $\phi$ through order $r+1$
when the input derivatives through spatial order $r+2$ satisfy that
envelope. The constants may depend on the derivative order.
Periodicity and stationarity, including either property on a time tail,
are preserved on their respective intervals.

For computable $L>0$, effective mixed-derivative evaluation and effective
derivative bounds for $g$ on finite time slabs give effective evaluation
and such bounds for $Pg,Qg,\phi$. Corresponding supplied global input
bounds give effective global bounds in each assertion above.
\end{proposition}
\begin{proof}
For $k\ne0$ define
\[
 \widehat\phi(t,k)=-\frac{i\xi_k\cdot\widehat g(t,k)}{|\xi_k|^2},
 \qquad \widehat\phi(t,0)=0.
\]
The nonzero-frequency symbols of $Q$ and $P$ are complementary orthogonal
projections. Their $H^s$ operator norms are at most one, while the
multiplier for $\phi$ has norm at most $d_L$ from $H^s$ to $H^{s+1}$.
For $|\beta|\le r$, Cauchy--Schwarz bounds the absolute sum of the
Fourier series for $\partial_x^\beta G$ by
\[
 \left(\sum_{k\in\Z^3}(1+|\xi_k|^2)^{-2}\right)^{1/2}
 \|G\|_{H^{r+2}}.
\]
The shell $|k|_\infty=n\ge1$ contains $24n^2+2\le26n^2$ points.
Since $\sum_{n\ge1}n^{-2}\le2$, the factor is at most $S_L$.
For the tail $|k|_\infty>N\ge1$, it is at most
\begin{equation}\label{eq:projection-tail}
 \left(26(L/(2\pi))^4/N\right)^{1/2}.
\end{equation}
Applying these estimates to the multipliers proves
\eqref{eq:projection-estimate}, including the estimate for $\phi$ with
one additional spatial derivative. On compact time slabs these tails
converge uniformly for every mixed derivative. Thus the series define
smooth fields and may be differentiated term by term. The Fourier
identities give the Poisson equation, uniqueness of its mean-zero
solution, and $\operatorname{div}Pg=0$. Also
$-\nabla(p-\phi)+Pg=-\nabla p+g$.

For every integer $n\ge0$, the multinomial expansion and Parseval give
\begin{equation}\label{eq:projection-derivative-bounds}
 \|G\|_{H^n}^2
 \le\sum_{|\beta|\le n}
       \frac{n!}{(n-|\beta|)!\,\beta!}
       \|\partial_x^\beta G\|_\infty^2.
\end{equation}
Use $n=r+2$ and $G=\partial_t^h g$. This finite sum proves the uniform
bounds and qualitative decay assertions. If its input suprema are
bounded by $C_{h,\beta}a_h(t)$, it gives the same envelope $a_h$ for
the output, with the constants specified by
\eqref{eq:projection-estimate}--\eqref{eq:projection-derivative-bounds}.
Integrating the squared estimates proves the temporal $L^2$ assertion
from the corresponding input norms. The operators act at each time,
so they preserve the stated temporal symmetries.

For effective evaluation, the supplied derivative bounds and
\eqref{eq:projection-tail} give a computable Fourier truncation on each
finite time slab. The finitely many coefficient integrals are computed
by Riemann sums. After scaling the torus to $[0,1]^3$, the integrand
for $\widehat G(k)$ has Lipschitz bound
\[
 L\sqrt3\max_j\|\partial_{x_j}G\|_\infty
       +2\pi|k|\|G\|_\infty.
\]
This gives an effective quadrature error and hence evaluation of every
mixed derivative of all three fields. The displayed finite sums also
compute output bounds from the supplied input bounds.
\end{proof}
The estimates apply to the pressure shift $\phi$; in the constructions
with original pressure zero, the projected pressure is $-\phi$.
Projection generally changes pressure. It supplies no compact-support
conclusion on Euclidean space and is distinct from the direct zero-pressure
solenoidal shear construction.

\paragraph{Effective flat profiles.}
For later cutoff formulas put
\begin{equation}\label{eq:p2-step}
 \rho(s)=\begin{cases}e^{-1/s},&s>0,\\0,&s\le0,\end{cases}
 \qquad \sigma(s)=\frac{\rho(s)}{\rho(s)+\rho(1-s)}.
\end{equation}
Positive-side derivatives of $\rho$ are polynomials in $1/s$ times
$e^{-1/s}$ and extend flatly at zero. The estimate
$v^me^{-v}\le(m+k)!v^{-k}$ gives effective error bounds near the seam,
and $\rho(s)+\rho(1-s)\ge e^{-2}$. Products, translations, rescalings,
and derivatives therefore yield effective cutoffs with all requested
derivative bounds. Periodization uses zero collars; at approximate real
arguments a finite superset of possible translates is evaluated. No real
equality test at a seam or execution of the encoded computation is used.
